% CS455-169
\section{Deployment, timeline and cost}\label{sec:deployment}

\subsection{Deployment}
The application runs on Render.com's Free tier, as a Docker web service with a Render managed PostgreSQL
database. Pushes to \texttt{master} deploy automatically, and migrations run when the container starts.
The health endpoint is \url{https://campus-ridepooling-automation.onrender.com/api/health}. The steps,
configuration and verification are in the
\repolink{Deliverables/D3/Deployment_Documentation.md}{deployment documentation}.

Deliverable 1 planned GCP (requirement NFR-RD-18). The team chose Render on 5 November because it builds
the repository's Docker image directly from GitHub, provides managed PostgreSQL and runs the custom server
that Socket.IO needs. The submitted requirement still says GCP, and the change is recorded in section 7 of
the deployment documentation rather than by rewriting Deliverable 1.

The Free tier has limits that matter for the demonstration. The service sleeps after 15 minutes without
traffic and takes about a minute to wake. The database expires on 5 December 2026 and has no backups.

\subsection{Timeline}
The \repolink{Deliverables/D3/project-timeline-gantt.md}{timeline and Gantt chart} (\jira{CS455-162}) covers
the project from the proposal to the final demonstration, and Figure~\ref{fig:gantt} on the last page shows it. Each work bar
runs from the day the pull requests for that work opened to the day they merged. Deliverable 1 work ran from
23 to 26 September, implementation from 2 to 15 October, and Deliverable 3 testing, reporting and deployment
from 28 October to 5 November.

\subsection{Cost}
The \repolink{Deliverables/D3/Costing_Documentation/project_costing_report.pdf}{costing report}
(\jira{CS455-178}) estimates the monthly operating cost of the current setup and of two production sizes,
from prices checked on 6 November 2026 at a rate of 94 rupees to the dollar.
Table~\ref{tab:cost} summarises it. Totals include 18\,\% GST.

\begin{table}[h]
\centering\small
\caption{Estimated monthly operating cost.}
\label{tab:cost}
\begin{tabular}{@{}lllr@{}}
\toprule
Setup & Web service & Database & Monthly cost \\
\midrule
Current demonstration & Render Free & Render Free, 1\,GB & INR 0 \\
Pilot, about 1,000 students & Render Starter & Basic-256mb, 5\,GB & INR 2,917 (USD 31.03) \\
Full scale, about 5,000 students & Render Standard & Basic-1gb, 20\,GB & INR 13,954 (USD 148.44) \\
\bottomrule
\end{tabular}
\end{table}

Model inference is the largest cost at both production sizes: INR 1,034 of the INR 2,472 pre-tax pilot cost,
and INR 5,170 of INR 11,825 at full scale. That figure assumes 20,000 and 100,000 Gemini calls a month, at
1,000 input and 200 output tokens per call. The input figure comes from the requests the code really builds;
the output figure is not measured, so the costing report also gives the effect of doubling it. Salaries,
monitoring tools and bandwidth above the included allowance are excluded.

The current setup stays within its free limits. A single web service stays within Render's 750 free hours a
month, a freshly migrated database is about 11\,MB of its 1\,GB limit, and the database expires after the
presentation window. GitHub Actions ran 468 workflows, about 1,040 elapsed minutes, at no charge, because
standard runners are free for a public repository. Section 7 of the costing report compares estimated and
actual usage and lists the production figures that were not collected.

Section 9 of the costing report gives an AI-assisted estimate of development time, made after the work: 181
member-days from the size of what was built, against 61 member-days with a commit. Commit days undercount
effort and no timesheets were kept, so the true figure lies somewhere between the two.
